\documentclass[a4paper,11pt]{article}

\usepackage{ascmac}
\usepackage{booktabs}
\usepackage[margin=25mm]{geometry}
\usepackage{fontspec}
\usepackage{xeCJK}
\setmainfont{Times New Roman}
\setCJKmainfont{Yu Mincho}
\setCJKsansfont{Yu Gothic}
\setCJKmonofont{MS Gothic}
\XeTeXlinebreaklocale "ja"
\XeTeXlinebreakskip=0pt plus 1pt
\setlength{\parindent}{1em}
\renewcommand{\today}{\number\year 年\number\month 月\number\day 日}

\title{Sentence-BERTを用いた用途・予算別PC構成推薦システム\\付録資料}
\author{学籍番号 24026111，氏名 佐藤 信晃}
\date{\today}

\begin{document}
\maketitle

\section{LLMの利用記録}
本課題では，LLMを以下の目的で利用した。
\begin{itemize}
  \item 課題文の要件整理と，既存プロトタイプとの差分確認
  \item 文字化けおよびPython構文エラーの修正
  \item 入力検証，キャッシュ，性能評価処理の実装支援
  \item レポート本文の構成と表現の推敲
\end{itemize}

主な指示は「NLPreport.mdの要件に沿って，既存のPCパーツ推薦
プロトタイプを完成させること」であった。
LLMから得た提案はそのまま採用せず，Pythonの構文検査，
依存ライブラリのインポート，Sentence-BERTによる実検索，
10件の一括評価を実行して確認した。

\subsection{対話記録}
以下では，実際にLLMへ入力した指示の内容を時系列で示す。
対話時の短い口語表現は，意図を変えない範囲でレポート用の文章に要約した。
したがって，以下の記述は逐語的な対話ログではなく，指示内容の要約である。
各指示に対する応答の要点，実際に採用した内容，確認方法も併記する。

\subsubsection*{対話1：課題要件と既存実装の確認}
\noindent
\textbf{ユーザの指示内容（要約）}\\
「\verb|NLPreport.md|に記載された課題要件を確認し，
既存のPCパーツ推薦システムを課題に適合するプロトタイプへ修正する。」

\noindent
\textbf{LLMの応答・作業}\\
課題文，既存の\verb|app.py|，\verb|pc_builds.csv|，
\verb|requirements.txt|を読み，課題がSentence-BERTまたはgensimを
必ず利用するNLPプロトタイプの設計・実装・評価であると整理した。
既存コードは日本語の文字化けに加え，文字列の引用符が途中で失われており，
Pythonとして実行できない状態であると指摘した。

\noindent
\textbf{採用内容と確認}\\
PC構成推薦という題材は課題条件を満たすため維持し，
文字化けしたUI文言と構文を修正した。
\verb|python -m py_compile app.py|による構文検査と，
Streamlit，pandas，sentence-transformers，torchのインポート確認を行った。

\subsubsection*{対話2：システムのゴールの明確化}
\noindent
\textbf{ユーザの質問内容（要約）}\\
「PC構成推薦システムの対象ユーザ，入出力，Sentence-BERTの利用箇所，
および性能評価の方法を整理する。」

\noindent
\textbf{LLMの応答}\\
利用者が自然文でPC用途と予算を入力すると，Sentence-BERTで用途の意味を比較し，
予算内のPC構成を推薦するシステムを完成形として説明した。
また，プロトタイプの作成だけでなく，Top-1・Top-3正解率による評価と，
指定形式の本文PDFおよび付録PDFが必要であると整理した。

\noindent
\textbf{採用内容}\\
この説明を基に，実装では「構成を検索」と「性能評価」の二つの画面を用意し，
レポートでは課題設定，技術概要，設計，実装，考察の順で説明することにした。

\subsubsection*{対話3：推薦処理と評価処理の改善}
\noindent
\textbf{ユーザの指示内容（要約）}\\
「Sentence-BERTを利用した推薦処理を実装し，プロトタイプの性能を
比較可能な指標で評価できるようにする。」

\noindent
\textbf{LLMの提案}\\
検索のたびに全候補をベクトル化すると処理が重複するため，
モデル，CSVデータ，全候補の埋め込みをキャッシュする構成が提案された。
また，空の入力や予算内候補がない場合の処理，
登録文をそのまま使わない10件の言い換え評価が追加された。

\noindent
\textbf{採用内容と確認}\\
\verb|st.cache_resource|でモデルとコーパス埋め込みを，
\verb|st.cache_data|でCSVをキャッシュした。
Sentence-BERTモデルを実際に読み込み，
「WQHDでゲームを快適に遊びたい」という入力に対して
「WQHDゲーム」が類似度0.883で1位になることを確認した。
さらに10件の一括評価を実行し，Top-1，Top-3ともに10/10件となった。

\subsubsection*{対話4：レポートテンプレートへの適用}
\noindent
\textbf{ユーザの指示内容（要約）}\\
「\verb|NLP_template_UTF8.tex|をレポートのテンプレートとして使用し，
実装内容と評価結果をレポート用の文章にまとめる。」

\noindent
\textbf{LLMの応答・作業}\\
テンプレートの指定section名を変更しないという課題要件を確認し，
「課題設定」「技術概要」「システムの設計」
「プロトタイプの実装」「最終的なシステムの考察」の五章で本文を作成した。
Sentence-BERTの仕組み，コサイン類似度の数式，処理フロー，
関数構成，評価方法，制約と改善案を追加した。

\noindent
\textbf{確認}\\
参考文献にはSentence-BERTの原論文，使用モデルのModel Card，
Sentence Transformers公式文書を記載した。
ソースコード全文を本文に載せず，推薦処理の中核部分だけを掲載した。

\subsubsection*{対話5：LaTeXからPDFへの変換}
\noindent
\textbf{ユーザの指示内容（要約）}\\
「作成したLaTeX文書をコンパイルし，提出内容を確認できるPDFを生成する。」

\noindent
\textbf{LLMの応答・作業}\\
実行環境にpLaTeXがなかったため，Tectonicをプロジェクト内へ導入した。
元テンプレートの\verb|jsarticle|と\verb|dvipdfmx|指定は
TectonicのXeTeX処理に直接対応しなかったため，
本文内容を維持したまま\verb|article|，\verb|xeCJK|，
Windowsの日本語フォントを使う設定へ変更した。
日付はテンプレートと同様に\verb|\date{\today}|を使い，
「年・月・日」の日本語表記になるよう再定義した。

\noindent
\textbf{確認}\\
コンパイル終了コード，PDFヘッダ，ファイル終端を確認し，
正常なPDFファイルが生成されたことを確認した。

\subsubsection*{対話6：実行画面と結果の追加}
\noindent
\textbf{ユーザの指示内容（要約）}\\
「プロトタイプの実行画面をレポートへ掲載し，表示された推薦結果を
レポート用の文章と表で説明する。」

\noindent
\textbf{LLMの応答・作業}\\
利用者が用意した二枚の実行画面を本文の実装章へ掲載した。
画面に表示された入力条件，三つの推薦候補，類似度，価格，予算残額を
表として整理した。
単に画像を貼るだけでなく，第1候補が入力用途と対応した理由，
第2候補が低予算eスポーツとなったことから分かる意味検索の限界を考察した。

\noindent
\textbf{確認}\\
画像をレポート用の英数字ファイル名で\verb|figures|ディレクトリへ保存し，
本文から図番号で参照した後，PDFを再コンパイルした。

\subsubsection*{対話7：本文PDFと付録PDFの分離}
\noindent
\textbf{ユーザの指示内容（要約）}\\
「課題の提出形式に従い，レポート本文と付録資料を別々のPDFファイルに分離する。」

\noindent
\textbf{LLMの応答・作業}\\
課題文では本文PDFと付録PDFを別ファイルで提出すると指定されていたため，
本文末尾に含めていたLLM利用記録，ソースコード，評価結果の生データを削除し，
新しい付録用TeXファイルへ移動した。

\noindent
\textbf{確認}\\
本文PDFと付録PDFを個別にコンパイルし，
本文側には指定された五章と参考文献だけが，
付録側には補助資料だけが含まれることを確認した。

\subsection{LLM出力の扱い}
LLMの提案には，実装環境に存在しないコマンドや，
LaTeX処理系によって利用できない設定が含まれる可能性がある。
そのため，生成されたコードや文章を無条件には採用せず，
エラーメッセージを確認して修正し，実行結果を根拠として採否を判断した。
特に推薦精度100\%という結果については，評価数が10件と少なく，
開発者が作成したデータであるため，一般的な精度を保証しないことを
本文の考察に明記した。

\section{ソースコード}
作成したソースコードは次のファイルから構成される。
\begin{itemize}
  \item \verb|app.py|：Streamlit画面，推薦処理，評価処理
  \item \verb|pc_builds.csv|：26件のPC構成コーパス
  \item \verb|requirements.txt|：実行に必要なライブラリ
\end{itemize}

実行コマンドを以下に示す。
\begin{screen}
\begin{verbatim}
.\venv\Scripts\python.exe -m streamlit run app.py \
  --server.fileWatcherType none
\end{verbatim}
\end{screen}

\subsection{推薦処理の主要部分}
\begin{screen}
\begin{verbatim}
def recommend(query: str, budget: int, top_k: int = 3):
    query = query.strip()
    if not query or budget <= 0 or top_k <= 0:
        return []

    all_rows = load_data()
    eligible_indices = [
        index
        for index, row in enumerate(all_rows)
        if int(row["price_yen"]) <= budget
    ]
    if not eligible_indices:
        return []

    model = load_model()
    corpus_embeddings = load_corpus_embeddings()
    query_embedding = model.encode(
        query,
        convert_to_tensor=True,
        normalize_embeddings=True,
    )
    eligible_embeddings = corpus_embeddings[eligible_indices]
    scores = util.cos_sim(
        query_embedding, eligible_embeddings
    )[0]
    top = scores.topk(
        k=min(top_k, len(eligible_indices))
    )
\end{verbatim}
\end{screen}

\section{評価結果の生データ}
\begin{center}
\begin{tabular}{rlll}
  \toprule
  No. & 期待カテゴリ & 1位カテゴリ & Top-3\\
  \midrule
  1 & 競技FPS & 競技FPS & 正解\\
  2 & ゲーム＋動画編集 & ゲーム＋動画編集 & 正解\\
  3 & 生成AI & 生成AI & 正解\\
  4 & WQHDゲーム & WQHDゲーム & 正解\\
  5 & 4Kゲーム & 4Kゲーム & 正解\\
  6 & ゲーム配信 & ゲーム配信 & 正解\\
  7 & 3D制作 & 3D制作 & 正解\\
  8 & 学生向け & 学生向け & 正解\\
  9 & シミュレーション & シミュレーション & 正解\\
  10 & VRゲーム & VRゲーム & 正解\\
  \bottomrule
\end{tabular}
\end{center}

\end{document}
